\subsection{Physics-Regularized Neural Fields from Experimental Thermocouple Measurements}
\label{sec:sci_ai_measured_heat_rod}
Physics-informed neural networks penalize differential-equation residuals alongside observed-data errors, but an inaccurate physical constraint can degrade predictive performance. To study this tradeoff using original physical measurements, this case study examines a laboratory dataset of eight temperature sensors on an electrically heated aluminum rod. The recorded temperatures were acquired every 10 seconds over 1,670 seconds and comprise 168 recorded time points per sensor, or 1,344 actual measurements. The laboratory used a 26-V, 250-mA configuration, with consecutive sensors separated by 12.7 mm. The original thermocouple log and experiment description are available from Sawaqed and collaborators in their \emph{Heat Conduction in a Metal Rod} dataset.

A tanh neural field $T_\theta(x,t)$ with two hidden layers of 32 units was fitted to six entire sensor traces (1,008 temperature measurements). One independent thermocouple trace (168 measurements) selected the regularization strength, and a different, previously unobserved thermocouple trace (168 measurements) was used for final evaluation. All measurements are recorded temperatures; no artificial temperature values were added. As a deliberately simplified regularizer, the model included the squared residual of $\partial_t T_\theta-\alpha\,\partial_{xx}T_\theta$ on nondimensionalized coordinates, where $\alpha$ is a fitted dimensionless parameter. This simplified homogeneous diffusion operator omits the detailed thermal boundary forcing and convective losses of the laboratory setup; the fitted coefficient must not be interpreted as measured aluminum thermal diffusivity.

Figure~\ref{fig:sci_ai_measured_heat_rod} presents the recorded temperature field, the original held-out thermocouple measurements and neural predictions, and the validation--test error comparison. The unregularized model ($\lambda=0$) achieved 0.1321\,\textdegree C validation RMSE and 0.1103\,\textdegree C held-out test RMSE. Adding the PDE residual with $\lambda=0.002$ increased the held-out RMSE to 0.1856\,\textdegree C, while $\lambda=0.02$ increased it to 0.4937\,\textdegree C. Since the unregularized model also had the lowest validation error, it was selected without consulting the held-out test. This negative result illustrates why physics-based penalties must be empirically validated against real observations rather than presumed beneficial.

\begin{figure}[htbp]
\centering
\includegraphics[width=0.96\linewidth]{figure_11_pinn_residual_field_error.pdf}
\caption{Actual laboratory measurements and measured performance of a physics-regularized neural field. (a) Eight recorded aluminum-rod thermocouple temperature histories; color represents measured temperature in degrees Celsius. (b) Measured temperatures from the CH6 thermocouple, which was excluded from both training and model selection, compared with three actually trained neural predictions. (c) Observed CH3 temperatures used exclusively for validation and the validation-selected model. (d) Validation and test RMSE for three prespecified regularization strengths, showing that the simplified homogeneous diffusion constraint degraded held-out accuracy. The plotted measurements are authentic recorded data; the curves are actual model predictions.}
\label{fig:sci_ai_measured_heat_rod}
\end{figure}
